\renewcommand{\thesection}{S\arabic{section}}
\setcounter{section}{1}
\section{Spectral separability bounds under common observation models}
\label{sec:Appendix_S2}

Throughout this appendix, we consider the simple hypotheses
\[
H_0:U=U_0,
\qquad
H_1:U=U_1,
\]
and let \(P_0\) and \(P_1\) denote the corresponding laws of the downstream observations.

For any decision rule \(\phi\), define the total testing error
\[
R(\phi)
=
P_0(\phi=H_1)
+
P_1(\phi=H_0).
\]
For two simple hypotheses,
\[
\inf_{\phi}R(\phi)
=
1-\lVert P_0-P_1\rVert_{\mathrm{TV}}.
\]
Pinsker's inequality therefore gives
\begin{equation}
\label{eq:S2-general-Pinsker}
\inf_{\phi}R(\phi)
\ge
1-
\sqrt{
\frac{1}{2}
D_{\mathrm{KL}}(P_0\Vert P_1)
}.
\end{equation}

We use the common spectral separability functional
\begin{equation}
\label{eq:S2-common-J}
J(U_0,U_1)
=
\int_{-\infty}^{\infty}
\frac{
|p\widehat G(f)|^2
}{
N(f)
}
\left|
\widehat U_0(f)-\widehat U_1(f)
\right|^2
\,df,
\end{equation}
where the normalization \(N(f)\) is specified by the observation
model below. Because each subsection treats a fixed observation model, we suppress the model subscript $\mathcal M$. 


\subsection{Gaussian observation noise}

\subsubsection{Setup and spectral representation}
Suppose that
\[
D(t)=p(U*g)(t)+\varepsilon(t),
\]
where \(\varepsilon\) is a mean-zero stationary Gaussian process with
covariance operator \(C_\varepsilon\). We assume that the Fourier
basis diagonalizes \(C_\varepsilon\), with spectral density
\(S_\varepsilon(f)>0\).

Because the two Gaussian laws have the same covariance and differ
only in their means,
\[
D_{\mathrm{KL}}(P_0\Vert P_1)
=
\frac12
\left\langle
\mu_0-\mu_1,
C_\varepsilon^{-1}(\mu_0-\mu_1)
\right\rangle.
\]
In the Fourier basis, this becomes
\[
D_{\mathrm{KL}}(P_0\Vert P_1)
=
\frac12
\int_{-\infty}^{\infty}
\frac{|p\widehat G(f)|^2}
{S_\varepsilon(f)}
\left|
\widehat U_0(f)-\widehat U_1(f)
\right|^2df.
\]


\subsubsection{KL divergence and separability}

Because the two Gaussian laws have the same covariance and differ
only in their means, their Kullback--Leibler divergence is
\[
D_{\mathrm{KL}}(P_0\Vert P_1)
=
\frac{1}{2}
\int_{-\infty}^{\infty}
\frac{
|p\widehat G(f)|^2
}{
S_\varepsilon(f)
}
\left|
\widehat U_0(f)-\widehat U_1(f)
\right|^2
\,df.
\]

Thus, with
\[
N(f)=S_\varepsilon(f),
\]
equation~\eqref{eq:S2-common-J} satisfies
\begin{equation}
\label{eq:S2-Gaussian-KL}
D_{\mathrm{KL}}(P_0\Vert P_1)
=
\frac{1}{2}J(U_0,U_1).
\end{equation}

Substituting equation~\eqref{eq:S2-Gaussian-KL} into
equation~\eqref{eq:S2-general-Pinsker} gives
\[
\inf_{\phi}R(\phi)
\ge
1-
\sqrt{
\frac{J(U_0,U_1)}{4}
}.
\]

Therefore, for any target total testing error
\(\alpha\in(0,1)\), if
\begin{equation}
\label{eq:S2-Gaussian-threshold}
J(U_0,U_1)
\le
4(1-\alpha)^2,
\end{equation}
then
\[
R(\phi)\ge\alpha.
\]
for every decision rule $\phi$. 


\subsection{Poisson observation noise}

\subsubsection{Setup and downstream mean difference}

Suppose that observations are conditionally independent across time
and satisfy
\[
D(t)\mid H_0
\sim
\operatorname{Pois}\!\left(m_0(t)\right),
\qquad
D(t)\mid H_1
\sim
\operatorname{Pois}\!\left(m_1(t)\right),
\]
where
\[
m_0(t)=p(U_0*g)(t),
\qquad
m_1(t)=p(U_1*g)(t).
\]

Assume that both mean processes are uniformly bounded:
\[
0<m_{\min}
\le
m_i(t)
\le
m_{\max}
<\infty,
\qquad
i\in\{0,1\},
\quad
\forall t.
\]

Define
\[
\Delta m(t)
=
m_1(t)-m_0(t)
=
p[(U_1-U_0)*g](t).
\]

Taking Fourier transforms gives
\[
\Delta\widehat m(f)
=
p\widehat G(f)
\left[
\widehat U_1(f)-\widehat U_0(f)
\right].
\]

Using the continuous-time Foureir convention, Parseval's identity therefore yields
\begin{equation}
\label{eq:S2-Poisson-Parseval}
\sum_t
\{\Delta m(t)\}^2
=
\int_{-\infty}^{\infty}
|p\widehat G(f)|^2
\left|
\widehat U_0(f)-\widehat U_1(f)
\right|^2
\,df.
\end{equation}

\subsubsection{KL divergence and quadratic bound}
For two scalar Poisson distributions,
\[
P_{\lambda_0}
=
\operatorname{Pois}(\lambda_0),
\qquad
P_{\lambda_1}
=
\operatorname{Pois}(\lambda_1),
\]
define
\[
d_{\mathrm{Pois}}(\lambda_0,\lambda_1)
:=
D_{\mathrm{KL}}
\!\left(
P_{\lambda_0}
\Vert
P_{\lambda_1}
\right).
\]

Its closed-form expression is
\[
d_{\mathrm{Pois}}(\lambda_0,\lambda_1)
=
\lambda_0
\log\!\left(
\frac{\lambda_0}{\lambda_1}
\right)
+
\lambda_1-\lambda_0.
\]

Equivalently, this divergence is the Bregman divergence generated by
\[
\psi_{\mathrm{Pois}}(m)
=
m\log m-m,
\]
whose second derivative is
\[
\psi_{\mathrm{Pois}}''(m)
=
\frac{1}{m}.
\]

Taylor's theorem therefore gives
\[
d_{\mathrm{Pois}}(\lambda_0,\lambda_1)
=
\frac{
(\lambda_1-\lambda_0)^2
}{
2\xi
}
\]
for some \(\xi\) between \(\lambda_0\) and \(\lambda_1\).

Since \(\xi\ge m_{\min}\),
\[
d_{\mathrm{Pois}}(\lambda_0,\lambda_1)
\le
\frac{
(\lambda_1-\lambda_0)^2
}{
2m_{\min}
}.
\]

Applying this bound at each time point and using conditional
independence gives
\begin{equation}
\label{eq:S2-Poisson-KL-upper}
D_{\mathrm{KL}}(P_0\Vert P_1)
\le
\frac{1}{2m_{\min}}
\sum_t
\{\Delta m(t)\}^2.
\end{equation}

\subsubsection{Spectral separability and testing error}

For Poisson observations, set
\[
N(f)=m_{\min}.
\]

Equations~\eqref{eq:S2-common-J} and
\eqref{eq:S2-Poisson-Parseval} then give
\[
J(U_0,U_1)
=
\frac{1}{m_{\min}}
\sum_t
\{\Delta m(t)\}^2.
\]

Consequently, equation~\eqref{eq:S2-Poisson-KL-upper} becomes
\begin{equation}
\label{eq:S2-Poisson-KL-J}
D_{\mathrm{KL}}(P_0\Vert P_1)
\le
\frac{1}{2}J(U_0,U_1).
\end{equation}

Combining equations~\eqref{eq:S2-general-Pinsker} and
\eqref{eq:S2-Poisson-KL-J} yields
\[
\inf_{\phi}R(\phi)
\ge
1-
\sqrt{
\frac{J(U_0,U_1)}{4}
}.
\]

Therefore, if
\begin{equation}
\label{eq:S2-Poisson-threshold}
J(U_0,U_1)
\le
4(1-\alpha)^2,
\end{equation}
then
\[
R(\phi)\ge\alpha.
\]
for every decision rule $\phi$.


\subsection{Negative-binomial observation noise}

\subsubsection{Setup and mean--variance structure}

Suppose that observations are conditionally independent across time
and satisfy
\[
D(t)\mid H_0
\sim
\operatorname{NB}\!\left(m_0(t),\kappa\right),
\qquad
D(t)\mid H_1
\sim
\operatorname{NB}\!\left(m_1(t),\kappa\right),
\]
where
\[
m_0(t)=p(U_0*g)(t),
\qquad
m_1(t)=p(U_1*g)(t),
\]
and the dispersion parameter \(\kappa>0\) is shared by the two
hypotheses.

Under this parameterization,
\[
\mathbb E[D(t)\mid H_i]
=
m_i(t),
\]
and
\[
\operatorname{Var}[D(t)\mid H_i]
=
V\!\left(m_i(t)\right),
\]
where the negative-binomial variance function is
\begin{equation}
\label{eq:S2-NB-variance-function}
V(m)
=
m+\frac{m^2}{\kappa}
=
m\left(1+\frac{m}{\kappa}\right).
\end{equation}

Assume
\[
0<m_{\min}
\le
m_i(t)
\le
m_{\max}
<\infty,
\qquad
i\in\{0,1\},
\quad
\forall t.
\]

Define
\[
V_{\min}
=
V(m_{\min})
=
m_{\min}
+
\frac{m_{\min}^2}{\kappa},
\]
and
\[
\Delta m(t)
=
m_1(t)-m_0(t)
=
p[(U_1-U_0)*g](t).
\]

\subsubsection{KL divergence and quadratic bound}

For a negative-binomial family with fixed dispersion parameter
\(\kappa\), the scalar KL divergence can be represented as the
Bregman divergence generated, up to affine terms, by
\[
\psi_{\mathrm{NB}}(m)
=
m\log m
-
(m+\kappa)\log(m+\kappa).
\]

Its second derivative is
\[
\psi_{\mathrm{NB}}''(m)
=
\frac{1}{m}
-
\frac{1}{m+\kappa}
=
\frac{1}{
m(1+m/\kappa)
}
=
\frac{1}{V(m)}.
\]

Taylor's theorem therefore gives
\begin{equation}
\label{eq:S2-NB-mean-value}
d_{\mathrm{NB}}(\lambda_0,\lambda_1)
=
\frac{
(\lambda_1-\lambda_0)^2
}{
2V(\xi)
}
\end{equation}
for some \(\xi\) between \(\lambda_0\) and \(\lambda_1\).

Because \(V(m)\) is increasing for \(m>0\), the bounded-mean
assumption implies
\[
V(\xi)
\ge
V(m_{\min})
=
V_{\min}.
\]
Hence,
\[
d_{\mathrm{NB}}(\lambda_0,\lambda_1)
\le
\frac{
(\lambda_1-\lambda_0)^2
}{
2V_{\min}
}.
\]

Applying this bound at every time point and using conditional
independence gives
\begin{equation}
\label{eq:S2-NB-KL-upper}
D_{\mathrm{KL}}(P_0\Vert P_1)
\le
\frac{1}{2V_{\min}}
\sum_t
\{\Delta m(t)\}^2.
\end{equation}

\subsubsection{Spectral separability and testing error}

Using the continuous-time Fourier convention, Parseval's identity gives
\begin{equation}
\label{eq:S2-NB-Parseval}
\sum_t
\{\Delta m(t)\}^2
=
\int_{-\infty}^{\infty}
|p\widehat G(f)|^2
\left|
\widehat U_0(f)-\widehat U_1(f)
\right|^2
\,df.
\end{equation}

For negative-binomial observations, set
\[
N(f)=V_{\min}.
\]

Equations~\eqref{eq:S2-common-J} and
\eqref{eq:S2-NB-Parseval} then give
\[
J(U_0,U_1)
=
\frac{1}{V_{\min}}
\sum_t
\{\Delta m(t)\}^2.
\]

Consequently, equation~\eqref{eq:S2-NB-KL-upper} becomes
\begin{equation}
\label{eq:S2-NB-KL-J}
D_{\mathrm{KL}}(P_0\Vert P_1)
\le
\frac{1}{2}J(U_0,U_1).
\end{equation}

Combining equations~\eqref{eq:S2-general-Pinsker} and
\eqref{eq:S2-NB-KL-J} gives
\[
\inf_{\phi}R(\phi)
\ge
1-
\sqrt{
\frac{J(U_0,U_1)}{4}
}.
\]

Therefore, if
\begin{equation}
\label{eq:S2-NB-threshold}
J(U_0,U_1)
\le
4(1-\alpha)^2,
\end{equation}
then
\[
R(\phi)\ge\alpha.
\]
for every decision rule $\phi$.